\pdfinfoomitdate=1
\pdftrailerid{}
\documentclass[12pt]{article}
\usepackage[letterpaper,left=18mm,right=18mm,top=22mm,bottom=20mm,headheight=15pt]{geometry}
\usepackage{tikz,fancyhdr}
\usepackage[T1]{fontenc}
\usepackage{lmodern}
\pagestyle{fancy}\fancyhf{}
\fancyhead[L]{\small Week 88 / Packing and kissing circles / Materials}
\fancyfoot[L]{\small Bellingham Math Circle / Week 88 / W88-M-v2}
\fancyfoot[R]{\small\thepage}
\renewcommand{\headrulewidth}{0.3pt}\renewcommand{\footrulewidth}{0pt}
\setlength{\parindent}{0pt}
% Disk sheets, added October 23, 2026 from the review card (fix 1). Coordinates are
% millimetres from the top-left corner of the paper, y down. Every circle is drawn with
% a 0.2 mm line whose centre is at the stated diameter, with a dot at its centre;
% editable/review/check_pdf.py measures them and the 100 mm bar in the PDF.
\newcommand{\sheet}[1]{\begin{tikzpicture}[remember picture,overlay,x=1mm,y=-1mm,line width=0.2mm]
\begin{scope}[shift={(current page.north west)}]#1
\draw (58,256)--(158,256);\draw (58,254)--(58,258);\draw (158,254)--(158,258);
\node[anchor=west,font=\small] at (161,256) {100 mm};
\end{scope}\end{tikzpicture}}
\newcommand{\disk}[3]{\draw (#1,#2) circle[radius=#3mm];\fill (#1,#2) circle[radius=0.4mm];}
\begin{document}
{\small Print at Actual Size on card stock and check that the bar below measures 100 mm. Cut just inside the line. A pair needs 13 disks of 40 mm and 5 of 80 mm; each table needs one 120 mm middle (page 3).\par}
\sheet{%
\foreach \x in {66,150}{\foreach \y in {82,166}{\disk{\x}{\y}{40}}}
\foreach \x in {42,86,130,174}{\disk{\x}{230}{20}}}
\newpage
\null
\sheet{%
\foreach \x in {42,86,130,174}{\foreach \y in {50,94,138,182,226}{\disk{\x}{\y}{20}}}}
\newpage
\null
\sheet{%
\disk{80}{90}{60}\node[font=\small] at (80,98) {120 mm middle};
\foreach \y in {50,94,138}{\disk{164}{\y}{20}}
\foreach \x in {66,150}{\disk{\x}{202}{40}}}
\end{document}
